\begin{pblock}{External validation: never worse, better where it matters}
Elevation against independent surveys of the flats themselves, and water
level against gauges never used in calibration. Centred RMSE, metres:

\vspace{2mm}
\begin{tabular}{@{}l l r r@{}}
\toprule
Site (truth) & Regime & uniform & \term{MAREA} \\
\midrule
Tagus interior (EMODnet, 286k px) & estuary, $\tau\!\approx\!45'$ & 0.279 & \emphnum{0.247} \\
Aveiro lagoon (LiDAR, 333k px)    & branched, $\tau\!:41\!-\!56'$ & 0.590 & \emphnum{0.580} \\
Vadehavet (EMODnet, 344k px)      & open flat                     & 0.275 & 0.275$^{*}$ \\
Villaviciosa (own RTK, 217 pts)   & short deep ria                & 0.211 & 0.211$^{*}$ \\
\bottomrule
\end{tabular}
\figcap{$^{*}$the operator \emph{declines} where the lag does not beat the
null or does not help — equality is the correct verdict there, verified,
not assumed. In Aveiro the gain concentrates where the lag grows:
$+24$ to $+53$\,mm per band.}

\vspace{3mm}
\pfig[0.75\linewidth]{fig_ladder}
\figcap{The Scheldt, blind: MAREA's lag gradient
(\emphnum{0.72\,min/km}) matches the tide gauges' 0.9, and at the
Terneuzen gauge — never used in calibration — its level sits 1\,cm from
the best any pure lag could do.}
\end{pblock}
